\subsection{Curvature quadrature domains}
\label{sec:curvature-quadrature-jets}

Retain $C,x,(L,h),r,z,e$ and $\nu=a(z)\,\mathrm{d}A$ from \cref{sec:compact-growing-jets}, and assume that the curvature is semipositive on $C$ and strictly positive on $C\setminus\{x\}$. We now allow integers $q=q_m\to\infty$ with $q=o(m)$. Write $B_{m,q}$ for the density of $H^0(C,L^m\otimes\mathfrak m_x^q)$ with the same norm. Put
\[
 \rho=\phi_{z\bar z},\qquad t=q/m,\qquad R=(t/r)^{1/(2r)}.
\]
The disc comparison above requires the weight error to be small throughout the disc. A larger range is obtained by absorbing the harmonic part of the error into the holomorphic section and locating the boundary on which the remaining potential has first-order contact. Its curvature quadrature identity determines that boundary.

\Needspace{18\baselineskip}
\begin{lemma}\label{lem:curvature-quadrature-domain}\index{Quadrature domain!curvature}
Let $Q(R,\zeta)$ be jointly smooth on a fixed disk of radius $b>1$,
with $Q(0,\zeta)=|\zeta|^{2r}$. Suppose every $Q_R$ is subharmonic
and $Q_{R,\zeta\bar\zeta}>0$ in a fixed collar of the unit circle.
For small $R$, the largest subharmonic minorant of $Q_R$ with a pole
$2r\log|\zeta|+O(1)$ at zero has exactly one noncontact component
$\Omega_R$. Its boundary is a radial graph near the unit circle. For
every fixed finite $k$ the normalized interior map $F_R$ and the
logarithm-subtracted harmonic part depend smoothly on $R$ in the
corresponding H\"older spaces, with $O(R)$ distance from their radial
values. All estimates used below require only finitely many such norms.
\end{lemma}
\begin{proof}
Write $\mathbb D$ for the unit disc. Fix $k\geq3$ and $0<\alpha<1$. Let $A^{k,\alpha}$ be the real Banach
space of functions holomorphic on $\mathbb D$, $C^{k,\alpha}$ on its closure,
whose value at zero is real. For $h\in A^{k,\alpha}$ near zero set $f(u)=u(1+h(u))$. Small $h$ gives a
univalent $f$, with $f(0)=0$ and $f'(0)>0$. On $|u|=1$ define
\[
 K(R,h)=f(u)(Q_R)_\zeta(f(u))-r,
 \qquad T(R,h)=(P_-K,\operatorname{Re}[K]_0),
\]
where $P_-$ projects onto strictly negative Fourier modes and $[\cdot]_0$
is the constant coefficient. The boundedness of Cauchy projections on
these H\"older spaces follows from the circle Hilbert kernel: subtract
the value at the evaluation point, split its integral at twice the
increment length, and use the H\"older difference bound on the inner
part and the differentiated kernel bound on the outer part. This gives
the $C^{0,\alpha}$ bound; commuting angular derivatives gives the
$C^{k,\alpha}$ bound.
The smooth superposition operator has no derivatives
of $h$. Thus $T$ is a smooth map to the negative-mode space times $\mathbb R$.
At $(R,h)=(0,0)$, $T=0$, and
\[
 D_hK(0,0)v=r^2(v+\bar v),\qquad
 D_hT(0,0)v=(r^2\overline{v-v(0)},\,2r^2v(0)).
\]
This is a bounded real-linear isomorphism. The Banach implicit-function
theorem therefore gives a unique normalized branch $f_R$, smoothly
depending on $R$ in the fixed space, with $\|f_R-\mathrm{id}\|=O(R)$.
Branches obtained in higher fixed spaces agree by lower-space uniqueness.
This supplies all finite spatial estimates and all finite parameter
expansions needed in the asymptotic argument. Each prescribed finite
order has its own sufficiently small parameter interval; no common
interval is inferred by intersecting infinitely many intervals.
The equation has not omitted an imaginary residue condition. The
function $K$ now extends holomorphically to $\mathbb D$, and $K(0)$ is purely
imaginary. Stokes' formula gives
\[
 \frac1{2\pi\mathrm{i}}\int_{|u|=1}(Q_R)_\zeta(f_R(u))f_R'(u)\,\mathrm{d}u
 =\frac1\pi\int_{f_R(\mathbb D)}(Q_R)_{\zeta\bar\zeta}\mathrm{d}A\in\mathbb R.
\]
The residue of $(r+K)f_R'/f_R$ is $r+K(0)$, so $K(0)$ is real as well,
and is zero. Consequently
\[
 S_R(\zeta)=\frac{r+K(f_R^{-1}(\zeta))}{\zeta}
 =r/\zeta+H_R(\zeta)
\]
has a holomorphic $H_R$ on $\Omega_R=f_R(\mathbb D)$ and equals $(Q_R)_\zeta$
on its boundary. Let
\[
 U_R=2r\log|\zeta|+2\operatorname{Re}\int_0^\zeta H_R(w)\,\mathrm{d}w+C_R.
\]
Its complex gradient agrees with that of $Q_R$ on the boundary.
The boundary difference is constant; choose $C_R$ to make it zero.
The difference and its gradient vanish there. Its inward normal second
derivative is $\Delta Q_R>0$, because the tangential second derivative
vanishes. Hence $Q_R-U_R>0$ in a uniform inner collar. At $R=0$,
\[
 U_0=1+2r\log|\zeta|,\qquad
 |\zeta|^{2r}-1-2r\log|\zeta|>0\quad(0<|\zeta|<1).
\]
Uniform convergence of the gap preserves its positive minimum on the remaining compact interior
sets; the logarithmic term controls a neighborhood of zero.

Patch $U_R=Q_R$ off $\Omega_R$. Value and gradient matching eliminate
boundary measures in its Laplacian. Its curvature inside is $r\delta_0$
and outside is $\ddc Q_R\geq0$, so it is an admissible minorant. If $v$
is any competitor, $v-U_R$ is subharmonic on $\Omega_R\setminus\{0\}$,
bounded above at zero and nonpositive on its boundary. Removable
subharmonic extension and the maximum principle give $v\leq U_R$.
Outside the hole this follows from $v\leq Q_R$. This identifies the
envelope and excludes extra noncontact components. The radial graph and
the harmonic-part estimates follow from the normalized map and its
parameter dependence. We henceforth denote this map by $F_R=f_R$.
\end{proof}

Taylor's integral remainder gives, for small signed $R$,
\[
 Q_R(\zeta)=R^{-2r}\phi(R\zeta)
 =|\zeta|^{2r}+R E(R,\zeta)
\]
with jointly smooth $E$. This applies the lemma to the actual weight.
Choose the harmonic primitive with $G_R(0)$ real, and write its envelope inside the hole as
\[
 U_R=2r\log|\zeta|+2\operatorname{Re} G_R(\zeta),\qquad
 G_R=1/2+O(R).
\]
Set $D_t=R\Omega_R$. In physical coordinates the obstacle potential is
$\phi_t(z)=R^{2r}U_R(z/R)$ inside $D_t$ and $\phi$ outside.
Replacing $\phi$ by $\phi_t$ in this chart and keeping $h$ elsewhere defines a global singular metric on $L$. Its potential is $C^1$ across the boundary, is bounded above by $\phi$, and has pole $2t\log|z|+O(1)$.
On the boundary $\phi_z=t/z+g_t'$ for a holomorphic $g_t$. Green's formula gives
\begin{equation}\label{eq:curvature-quadrature-identity}
 \int_{D_t}\rho(z)f(z)\mathrm{d}A=\pi t f(0)
\end{equation}
for every $f$ holomorphic near $\overline{D_t}$. The curvature is $t\delta_x$ inside and the original curvature outside. With $c_R$ a positive periodic function, the boundary is written as $\Gamma_t=\{R c_R(\theta)\mathrm{e}^{\mathrm{i}\theta}\}$, where $c_R-1=O(R)$ in each fixed norm used below.
